\documentclass[11pt]{article}
\usepackage[margin=1in]{geometry}
\usepackage{amsmath,amssymb,amsthm}
\usepackage{booktabs,array}
\usepackage{xurl}
\usepackage[hidelinks]{hyperref}
\usepackage{url}
\setlength{\emergencystretch}{2em}
\newcommand{\R}{\mathbb R}
\newcommand{\C}{\mathcal C}
\newcommand{\K}{\mathcal K}
\newcommand{\D}{\mathcal D}
\newcommand{\Sf}{\mathcal S}
\newcommand{\B}{\mathcal B}
\newcommand{\Pre}{\operatorname{Pre}}
\newtheorem{proposition}{Proposition}
\title{A joint-state terminal safety set under the constant-parameter target contract}
\author{}
\date{Prepared October 10, 2026}
\begin{document}
\maketitle

\paragraph{Status and scope.}
This is a design study, not an adopted controller or a new safety certificate.
The production baseline examined is commit
\texttt{423e3e885d63d3461f8567726f3a9a5946c757ec}: steering angle and braking
ratio remain the inputs, and the terminal law is $u=u^*+Ke$.
No vehicle campaign or MATLAB suite was rerun for this study.
The question is how to define a terminal set that expresses safety of the
ego--target interaction, uses the constant tangential acceleration and
sideslip contract, and admits a short invariance argument without making
nominal cruise a prerequisite for every safe state.

\paragraph{Recommendation.}
Use the joint viability kernel as the ideal mathematical object. Construct
a certified inner approximation by safe predecessors of a joint invariant
core, or by the safe flow of one joint feedback into that core. Retain the
existing cruise certificate as a useful local ingredient after separation
has been certified. With imperfect observations, apply the construction to
an information set of joint states, with one common causal policy for all
states consistent with the measurements. This does not require a high
dimensional grid, although finding a useful certified approximation is a
substantial remaining problem.

\section{What the current construction establishes}
The current ellipsoid $\Omega=\{x:e(x)^TPe(x)\leq1\}$ is a contraction
domain for the cruise feedback; its certification addresses the sampled
nonlinear ego model and the input and error boxes. Collision clearance is
established separately by checking the ego CLF tube against the target
forecast. Thus the implemented terminal set is already parameterized by
target information: with an exact target state $q$, it can be written as
$\Sf(q)$, a section of a joint set
\begin{equation}
\Sf_{\rm joint}=\{(x,q):x\in\Sf(q)\}.
\end{equation}
A parameterized set or a time-dependent set indexed by an exact forecast
can be equivalent to an explicitly augmented state representation. Merely
concatenating state vectors adds no safety theorem. The essential question
is whether the joint evolution preserves the defining safety conditions.
An ego-only set could also be safe if it is robust to every allowed target,
but that stronger requirement may make it very small or empty.

The baseline proves recursive feasibility conditionally: exact ego state
and declared plant (H1), the unchanged constant-parameter target forecast
(H2), and containment of the sampled nonlinear maps in the recorded
enclosure (H3). Section 9 of \path{controller/TERMINAL_SAFE_SET.md} analyzes
joint estimation and control; Section 10 states that the target part is
certainty-equivalent. Its numerical H3 is supported by samples and a margin,
not by an interval proof. This study does not upgrade those premises.

The present terminal policy is restricted to returning to the given path
trim at the configured speed. A safe same-lane following state can therefore
be excluded even when admissible braking and a lower speed would maintain
separation. Enlarging the ellipsoid changes this policy's certified domain;
it cannot by itself recover states requiring a different policy.

\section{Autonomous joint dynamics and exact target flow}
Let the ego state be $x=(p_E,\psi_E,v_x,v_y,r_E)$ and its controlled
nonlinear bicycle dynamics be $\dot x=f_E(x,u)$, with
$u=(\delta,\lambda_b)$ the existing steering angle and braking ratio.
Write the dynamic target state, omitting known dimensions, as
$q=(p_T,\psi_T,V_T,A_T,\beta_T)$. The target contract is
\begin{align}
\dot p_T&=V_T\begin{bmatrix}\cos(\psi_T+\beta_T)\\
                                  \sin(\psi_T+\beta_T)\end{bmatrix}, &
\dot\psi_T&=\frac{V_T\sin\beta_T}{l_{r,T}},\label{eq:target}\\
\dot V_T&=A_T, & \dot A_T&=0,\qquad\dot\beta_T=0.\nonumber
\end{align}
These equations are the inertial model in Sharma et al., Eq.~(2); its
ego-frame form is their Eq.~(17) \cite{sharma}. They agree with
\texttt{predictiveSafetyGeometry.predictTarget}.
The augmented state $z=(x,q)$ obeys an autonomous system
$\dot z=F(z,u)$: no independent future target input is introduced.

Let $k_T=\sin\beta_T/l_{r,T}$, $\chi_0=\psi_T+\beta_T$, and
\begin{equation}
\ell(t)=V_Tt+\tfrac12 A_Tt^2.
\end{equation}
The exact flow is
\begin{align}
V_T(t)&=V_T+A_Tt, & \psi_T(t)&=\psi_T+k_T\ell(t),\\
p_T(t)&=p_T+\ell(t)\operatorname{sinc}\!\left(\frac{k_T\ell(t)}2\right)
\begin{bmatrix}\cos(\chi_0+k_T\ell(t)/2)\\
                \sin(\chi_0+k_T\ell(t)/2)\end{bmatrix},\label{eq:flow}
\end{align}
where $\operatorname{sinc}(a)=\sin(a)/a$ with value 1 at zero.
Thus the spatial orbit is a line for $k_T=0$ and a circle otherwise;
acceleration determines progress along the orbit. This gives exact target
propagation, a consistent semigroup $\phi_{t+s}(q)=\phi_t(\phi_s(q))$,
and analytic geometric ingredients for tail certificates. It does not
make the nonlinear ego dynamics or the joint safe kernel elementary.

Unknown constants remain unknown constants. For example, if
$A_T\in[A_-,A_+]$, the admissible trajectories are parameterized by one
fixed $A_T$ throughout each trajectory. Replacing this with an adversary
free to pick a new acceleration at every instant enlarges the motion class
and discards valuable temporal correlation. The same applies to sideslip.
Target position, heading, speed and parameter uncertainty should be
propagated with their correlations whenever feasible.

\paragraph{Contract validity.}
The production predictor uses signed $V_T$, so a negative constant
acceleration continues through zero into reversal. Equation~\eqref{eq:flow}
is exact for that declared contract. A physical rule that stops the target
at zero speed is a different, hybrid model and requires a corresponding
terminal analysis. Positive constant acceleration likewise implies
unbounded speed if extrapolated indefinitely. A finite model-validity
interval needs a certified continuation contract; it cannot alone support
an infinite-horizon safety claim.

\section{A natural relative-state representation}
Let $R(\psi)$ be planar rotation,
$J=\left[\begin{smallmatrix}0&-1\\1&0\end{smallmatrix}\right]$,
$\rho=R(-\psi_E)(p_T-p_E)$ and $\eta=\psi_T-\psi_E$.
Then
\begin{align}
\dot\rho&=V_T\begin{bmatrix}\cos(\eta+\beta_T)\\\sin(\eta+\beta_T)\end{bmatrix}
       -\begin{bmatrix}v_x\\v_y\end{bmatrix}-r_EJ\rho,\label{eq:relative}\\
\dot\eta&=k_TV_T-r_E, & \dot V_T&=A_T, &
\dot A_T&=\dot\beta_T=0.\nonumber
\end{align}
Rectangle separation depends on $(\rho,\eta)$ and the known body shapes,
while road clearance depends on ego path coordinates and heading. For a
spatially homogeneous road, a useful representation is
\begin{equation}
\zeta=(d_E,\theta_E,v_x,v_y,r_E,\rho_x,\rho_y,\eta,V_T,A_T,\beta_T),
\end{equation}
where $\theta_E$ is ego body heading relative to the path tangent.
Along-path station must be retained when road width, curvature, endpoints
or other constraints depend on it. Coordinates must also preserve any
needed winding or branch information on circular paths. Removing global
coordinates without these conditions would change the problem.

This representation connects directly to the NRMM observer's relative
position and yaw states. Frozen $(A_T,\beta_T)$ define parameterized
leaves of the dynamics rather than extra controlled directions. The
target's analytic orbit offers structure for computation; it does not
eliminate all the degrees of freedom in joint viability.

\section{The ideal set: joint viability}
Let $\C$ be the set of joint configurations with positive rectangle
clearance, ego road containment and the required state constraints.
Use a strictly positive collision buffer when expressing clearance through
ordinary set distance. Input and mixed handling constraints belong to the
admissible controls and hold conditions. With known states and exact
dynamics, the ideal terminal set is
\begin{equation}
\boxed{\K=\{z_0:\exists\pi\text{ admissible and causal},\quad
                  z^\pi(t;z_0)\in\C\ \forall t\geq0\}.}\label{eq:kernel}
\end{equation}
Under standard well-posedness and control-selection conditions, this is
the largest controlled invariant subset of $\C$, the viability kernel
\cite{aubin}. Maximal here means set inclusion under the specified model,
constraints and information pattern; it does not mean fastest travel,
minimum intervention, or any particular ellipsoid volume.

For a sampled implementation, define $\Psi_\tau(z,u)$ as the declared
within-hold evolution with constant input $u$, for $0\leq\tau\leq h$.
The ego and target models used to define $\Psi$ must match the certification
model. The safe predecessor is
\begin{equation}
\Pre(S)=\{z:\exists u\text{ admissible over the hold},\quad
\Psi_\tau(z,u)\in\C\ \forall\tau\in[0,h],\quad\Psi_h(z,u)\in S\}.\label{eq:pre}
\end{equation}
The greatest fixed point of $S\mapsto\C\cap\Pre(S)$ is the kernel for
this held-control class. It can differ from the kernel with unrestricted
continuous-time feedback. Checking only the two hold endpoints does not
establish the intersample condition.

The kernel may be empty for some encounters, such as an unavoidable
oncoming vehicle in a narrow corridor with insufficient braking or lateral
authority. The contract determines the threat; it does not guarantee that
the ego has an escape. This is a property to report, not an excuse to
weaken the collision constraint.

\section{A constructive terminal set with a short proof}
Computing the full kernel of the Fiala bicycle and rotating rectangles
is not presently tractable through a generic dense grid. Instead choose
a joint core $\D\subset\C$ with a certified admissible feedback and safe
holds that return to $\D$. Define
\begin{equation}
\Sf_0=\D,\qquad\Sf_{m+1}=\Sf_m\cup\Pre(\Sf_m).\label{eq:expansion}
\end{equation}
Each $\Sf_m$ contains the joint states from which a feasible continuation
can reach the core within at most $m$ holds. The continuation is a witness
that must be preserved or replaced by another verified witness.

\begin{proposition}
Every $\Sf_m$ in \eqref{eq:expansion} is controlled invariant at sample
instants, with safety throughout the holds, provided the core and
predecessor conditions are sound.
\end{proposition}
\begin{proof}
For $m=0$, use the core feedback. Inductively a state already in $\Sf_m$
has a safe input with successor in $\Sf_m\subset\Sf_{m+1}$. A newly
admitted predecessor has a safe input with successor in $\Sf_m$ by
\eqref{eq:pre}. Thus every state in $\Sf_{m+1}$ has a safe hold and a
successor in $\Sf_{m+1}$. Repeat the argument at each sample.
\end{proof}
This also gives the terminal part of the MPC shift proof: shift the
accepted joint trajectory and append one of these feasible terminal
holds. Entire-plan safety additionally requires hard collision and road
constraints over the executed part of the plan. The baseline's soft
prefix collision deficits are recovery measures and do not by themselves
guarantee collision freedom within that prefix.

A simpler, more conservative specialization fixes one joint backup
feedback $\kappa_b(z)$ that preserves $\D$. For an integer $M$ its
safe-capture set is
\begin{equation}
\Sf_b^M=\{z:\text{all }M\text{ backup holds are safe and }
\Phi_b^M(z)\in\D\}.\label{eq:backup}
\end{equation}
Shifting those $M$ holds and adding a core hold proves invariance at
sample instants. If a separate core controller is used, the continuation
must include its certified transition, or use the safe-predecessor
construction. This is the backup-set principle; the standard continuous
flow formulation is given by Chen et al., Eq.~(9) and Theorem~1
\cite{backup}. Its application to this project's dynamics still needs
model, input and intersample certificates.

\paragraph{What must constitute the core.}
A cruise contraction domain is usable only with joint separation. One
candidate core combines that domain, road containment, and a certificate
that the entire future target orbit and timing avoid the future ego tube.
For a straight road a relative-motion tail calculation can certify this;
for a turning target, separation from its swept orbit is a sufficient
geometric condition. Circular ego paths need future winding and recurrence
handled explicitly. Another core could admit safe following at the target
speed, but that behavior needs its own admissibility and invariance proof.
No such general following core is implemented or certified by this study.

\paragraph{Completeness and computation.}
The sets in \eqref{eq:expansion} are inner approximations of the kernel.
States that stay safe forever without ever reaching $\D$ may be omitted,
even as $m$ grows. A maximum kernel computation therefore cannot be claimed
from finite safe capture alone. Larger $M$ also increases online work;
this study supplies no timing bound or performance prediction. A robust
nominal-trajectory tube or a certified implicit barrier can represent an
inner approximation without enumerating the full joint state space.

\section{Cruise is a preference within the viable controls}
For a chosen certified $\Sf$, define
\begin{equation}
\mathcal U_{\Sf}(z)=\{u:\text{the hold is admissible and safe, and }
\Psi_h(z,u)\in\Sf\}.
\end{equation}
A joint terminal feedback may then be selected by
\begin{equation}
\kappa_f(z)\in\underset{u\in\mathcal U_{\Sf}(z)}{\operatorname{argmin}}
\|u-\kappa_{\rm cruise}(x)\|_R^2.\label{eq:governor}
\end{equation}
Nonemptiness of the feasible control set comes from invariance, not from
the cruise objective. This one control rule tries to cruise whenever that
is compatible with continued safety. It may reduce speed or alter steering
when the interaction requires it. Its inputs remain steering angle and
braking ratio, without an inverse tire-force map or prescribed lane modes.
The nonlinear predecessor and rectangle constraints need not be convex;
\eqref{eq:governor} is not automatically a QP or a cheap online algorithm.
A smooth CBF is one possible representation of a certified set, not an
independent reason for feasibility.
The existing $(P,K)$ contraction and its tube apply to the original cruise
law on its certified domain. A governor that changes that input cannot
reuse the same contraction or tube without a new proof or an explicit
bound on the change. A joint backup and its capture set therefore need
their own certificate; reusing the cruise core does not bypass this step.

Safety and asymptotic cruise recovery are different theorems. An invariance
constraint alone does not ensure eventual cruise, progress or comfort.
The existing CLF can support recovery once the interaction permits it;
requiring fixed-speed CLF decay at every joint state would again exclude
otherwise viable following states. The PCBF literature likewise requires
an invariant terminal set and distinguishes set safety from equilibrium
stabilization \cite{pcbf}.

\section{A transparent one-dimensional example}
Consider aligned vehicles on a straight path with a constant-speed target
($A_T=\beta_T=0$). Let $g$ be the bumper clearance minus the required
buffer and $w=v_E-v_T$ the closing speed. Suppose the simplified ego
model can brake at $b>0$ until it matches an admissible target speed, then
maintain that speed. Define
\begin{equation}
B(g,w)=g-\frac{[w]_+^2}{2b}.\label{eq:braking}
\end{equation}
The set $B\geq0$ is invariant under braking while $w>0$ and matching the
target speed afterwards. Indeed $\dot g=-w$, $\dot w=-b$, so while closing,
$\dot B=-w-(w/b)(-b)=0$. Once $w=0$, both $g$ and $B$ remain constant.
For $w<0$, maintaining the speeds increases the gap. A state with
$g<[w]_+^2/(2b)$ cannot be saved by longitudinal control bounded below by
$-b$ in this simplified same-lane model, since maximum braking minimizes
the accumulated closing displacement. Thus this example has a direct
viability interpretation, and is a relative-motion version of the standard
stopping-distance barrier \cite{backup}.

The terminal safety criterion depends on both vehicles, and a safe cruise
speed can be the target's speed rather than the mission's preferred speed.
Equation~\eqref{eq:braking} is an illustration, not a certificate for the
production bicycle: tire force, yaw, slip, sampling, speed limits, target
acceleration and model errors must be included before using it there.

\section{Imperfect observations: an information-set kernel}
Let $\B_k$ contain all joint states and fixed target parameters consistent
with the history of measurements, bounded errors and the motion contract.
The correct output-feedback requirement is
\begin{equation}
\exists\ \text{one causal observation-based policy}\quad
\forall z_k\in\B_k:\quad \text{the resulting joint trajectory is safe}.
\end{equation}
It is stronger than $\forall z\in\B_k,\ \exists\kappa_z$.
The latter permits different controls for indistinguishable states and
does not give an executable policy. In particular, $\B\subset\K$ alone
does not establish robust output-feedback viability.

With one common held input $u$, predict the set $\overline\B^+(\B,u)$
and intersect it with the set consistent with the next possible
measurement $y$:
\begin{equation}
\B^+(\B,u,y)=\overline\B^+(\B,u)\cap\mathcal M(y).
\end{equation}
The prediction includes every allowed ego disturbance, and holds the
unknown $A_T,\beta_T$ fixed on each trajectory. Define the information
predecessor of a family of information sets $\mathfrak S$ by
\begin{align}
\Pre_I(\mathfrak S)=\{\B:\ &\exists\text{ one admissible }u,
\text{ all within-hold joint trajectories are in }\C,\nonumber\\
&\B^+(\B,u,y)\in\mathfrak S\text{ for every consistent }y\}.\label{eq:info}
\end{align}
The greatest fixed point is the information-state safety kernel. The
core-and-predecessor induction applies here as well. Product sets for ego
and target uncertainty are acceptable sound outer bounds, but can lose
correlations from shared sensors and coordinates.

Measurement intersection gives
$\B_{k+1}\subseteq\overline\B^+(\B_k,u_k)$: containment compares the
sets at the same successor time. It does not assert
$\B_{k+1}\subseteq\B_k$ in different-time coordinates. Future target
predictions generated from this update can be nested at matching absolute
times. This is a convenient sufficient route to reusing an earlier
certificate, not a universal necessity for every robust control method.
Methods that explicitly accommodate future information changes can also
establish safety. Replacing one point forecast by a freshly estimated one
without these consistency conditions does not preserve a shift proof.
The need to account for measurement errors is independently established
by the measurement-robust CBF framework \cite{measurement}.

\section{Perception exit is not physical non-reencounter}
The baseline defines a finite encounter with an absorbing exit when the
target leaves the perception range. This is a stated task convention;
it is not a proof of physical all-time separation from that target.
Under the repository's signed-speed target contract, take an ego at
8~m/s and a straight target initially 45~m ahead, with initial speed
18~m/s and constant acceleration $-2$~m/s$^2$. The reference-point gap is
\begin{equation}
d(t)=45+10t-t^2.
\end{equation}
It exceeds 50~m after 0.528~s, returns below 50~m after 9.472~s, and the
reference points meet at 13.367~s. Rectangle contact occurs earlier.
This is a counterexample to treating range exit alone as a physical
infinite-horizon safety certificate. It is not a contradiction of the
baseline's explicitly finite-encounter theorem.

For physical safety, remove the target from the joint state only after a
certificate of continued separation under the declared contract, or retain
it in a valid uncertainty model. The collision-cone certificate is useful
when it actually checks all later motion for the prescribed continuation.
Merely receding at the present instant is insufficient. Newly entering,
unmodeled vehicles need additional perception and initial-viability
assumptions; a one-target terminal set cannot cover arbitrary unseen
traffic.

\section{Derivation checks and remaining obligations}
Independent arithmetic checks, seed 20261010, compared the analytic target
flow with a separate RK4 integration of \eqref{eq:target}, and checked
the rotation derivative in \eqref{eq:relative} by central finite
differences. Straight and curved cases, near-zero sideslip, signed speeds
and passage through zero speed were included. They are numerical
consistency checks, not exhaustive or interval certificates.

\begin{center}
\begin{tabular}{@{}p{0.54\linewidth}rp{0.25\linewidth}@{}}
\toprule
Check & Cases & Largest absolute discrepancy\\
\midrule
Target-flow semigroup & 64 & $2.85\,10^{-14}$\\
Target flow versus independent RK4, 4096 steps over 4~s & 14 & $7.52\,10^{-12}$\\
Relative position/yaw vector field, finite-difference step $10^{-5}$~s & 64 & $7.57\,10^{-9}$\\
Scalar braking barrier along brake-to-match trajectories & 100 & $2.25\,10^{-14}$\\
\bottomrule
\end{tabular}
\end{center}
Position, angle and velocity component errors in the target-flow checks
use their respective native SI units; the maximum listed is a numerical
component maximum, not a dimensionally meaningful norm across unlike units.
Source and exact results are stored with this study; the computation does
not measure collision-avoidance success rates or real-time feasibility.
The reproducible command from the repository root is
\begin{quote}\small
\path{python scripts/checkJointTerminalSetDerivations.py --output}\\
\path{report/JOINT_TERMINAL_SET_DESIGN_20261010/derivation-checks.json}
\end{quote}

Before implementation, a useful nonempty joint core and an admissible
continuation policy must be certified on the declared bicycle model.
The hold clearance and target-tail bounds must be sound; sampling with a
margin alone should be identified as a numerical hypothesis. For noisy
operation, common-policy uncertainty propagation and a compatible
information update are also required. An efficient inner approximation
must be evaluated on the same exact and noisy campaigns; no improvement
is predicted here. Existing prescribed lane backups, force-level input
maps and removed target-estimation modules are not reinstated by this
proposal. The research conclusion is a proposed mathematical architecture,
not an implementation decision by the user.

\paragraph{Source verification.}
The primary papers below were checked against their arXiv preprints or
publisher pages. Sharma's Eqs.~(2), (16), (17) were read from the local
reference PDF. The local Zotero API was probed but was not running; no
library items were modified. The joint predecessor and information-set
formulation here are an application and synthesis of invariant-set ideas
to this repository, not a claim of a new general theorem.

\begin{thebibliography}{9}
\bibitem{sharma}
G. Sharma, H. Alai, and R. Rajamani,
``Simultaneous ego-vehicle state estimation and vehicle trajectory tracking
using a multistage high gain observer,''
\emph{Transportation Research Part C}, vol. 182, 105411, 2026.
\href{https://doi.org/10.1016/j.trc.2025.105411}{doi:10.1016/j.trc.2025.105411}.
\bibitem{aubin}
J.-P. Aubin, ``A Survey of Viability Theory,''
\emph{SIAM Journal on Control and Optimization}, 1990.
\href{https://doi.org/10.1137/0328044}{doi:10.1137/0328044}.
\bibitem{backup}
Y. Chen, M. Jankovic, M. Santillo, and A. D. Ames,
``Backup Control Barrier Functions: Formulation and Comparative Study,''
2021. \url{https://arxiv.org/abs/2104.11332}.
\bibitem{pcbf}
J. Huang, H. Wang, K. Margellos, and P. Goulart,
``Predictive Control Barrier Functions: Bridging model predictive control
and control barrier functions,'' arXiv version 2, 2025.
\url{https://arxiv.org/abs/2502.08400}.
\bibitem{measurement}
R. K. Cosner, A. W. Singletary, A. J. Taylor, T. G. Molnar,
K. L. Bouman, and A. D. Ames,
``Measurement-Robust Control Barrier Functions: Certainty in Safety with
Uncertainty in State,'' 2021.
\url{https://arxiv.org/abs/2104.14030}.
\end{thebibliography}
\end{document}
